(1) 无穷远处能量和角动量分别为 $E=mv_\infty^2/2$、$L=mv_\infty b$。由于中心力守恒，
\[
\frac m2\dot r^2+\frac{L^2}{2mr^2}+\frac\kappa r=E.
\]
在最近距离处 $\dot r=0$，方程 $2Er^2-2\kappa r-L^2/m=0$ 仅有一个正根：
\[
r_{\min}=\frac{\kappa+\sqrt{\kappa^2+m^2v_\infty^4b^2}}{mv_\infty^2}.
\]
$L\ne0$ 时有效势在 $r\to0$ 趋于正无穷，因此粒子不会撞击中心。

(2) 令 $u=1/r$，Binet 方程为 $u''+u=-m\kappa/L^2$。选近心点为极角原点，解为
\[
u(\phi)=\frac{m\kappa}{L^2}[e\cos\phi-1],\qquad
e=\sqrt{1+\frac{2EL^2}{m\kappa^2}}
=\sqrt{1+\frac{m^2v_\infty^4b^2}{\kappa^2}}>1.
\]
轨道两端 $r\to\infty$ 对应 $\cos\phi_\infty=1/e$。入射与出射渐近线形成的散射偏转角大小为
\[
\Theta=\pi-2\phi_\infty
=2\arcsin(1/e)
=2\arctan\frac{\kappa}{mv_\infty^2b}.
\]
该角度严格介于零与 $\pi$ 之间，是双曲轨道的渐近散射角。

(3) 冲量参数在 $[b,b+db]$ 对应入射面积 $d\sigma=2\pi b\,|db|$，散射到极角 $[\Theta,\Theta+d\Theta]$ 对应立体角 $d\Omega=2\pi\sin\Theta\,|d\Theta|$。反解
\[
b=\frac{\kappa}{mv_\infty^2}\cot\frac\Theta2,\qquad
\left|\frac{db}{d\Theta}\right|=
\frac{\kappa}{2mv_\infty^2}\csc^2\frac\Theta2.
\]
因此
\[
\frac{d\sigma}{d\Omega}=\frac{b}{\sin\Theta}\left|\frac{db}{d\Theta}\right|
=\frac{\kappa^2}{4m^2v_\infty^4}\csc^4\frac\Theta2
=\frac{\kappa^2}{16E^2}\csc^4\frac\Theta2.
\]
这一量的量纲为面积，可从 $\kappa/E$ 是长度来检查。

(4) $b\to0^+$ 给 $\Theta\to\pi$、$r_{\min}\to2\kappa/(mv_\infty^2)$；$b\to\infty$ 或固定 $b$ 令 $\kappa\to0^+$ 时 $\Theta\to0$。前向小角处 $\sin(\Theta/2)\sim\Theta/2$，于是 $d\sigma/d\Omega\propto\Theta^{-4}$，其对前向小角区域的积分发散。实际屏蔽长度与探测器角分辨率会引入额外截断；理想无屏蔽的库仑总散射截面不是有限常数。
